%!TEX root = ../template.tex
%% ====================================================================
%% Capitulo 3 -- Revisao sistematica do ambito
%% Esqueleto gerado 2026-08-13 a partir do indice v1
%% (Mind Palace: salas/tese/2026-08-13-indice-dissertacao-v1.md)
%% Cada '>>' descreve o que a seccao deve conter. Apagar ao escrever.
%% ====================================================================

\chapter{Revisão sistemática do âmbito}
\label{cha:revisao}

% [REV-2026-09-13][C3-01][P2][CONTEUDO] Fechar a revisão com estado verificável
% Indicar ID/URL e data do registo, perguntas, critérios de elegibilidade,
% datas de pesquisa e emendas. PRISMA-ScR orienta o relato; não certifica
% o método nem substitui um protocolo. Se só existe piloto de extração,
% assumi-lo aqui e ajustar Introdução/conclusões; não apresentar uma síntese
% completa por antecipação.
% Fonte verificada: https://www.prisma-statement.org/scoping
% [/REV-2026-09-13]
\section{Protocolo e pré-registo}
\label{sec:protocolo-sr}

%% >> Scoping review registada no OSF segundo PRISMA-ScR, com SQ1--SQ4 e a estrutura de dois corpora.


% [REV-2026-09-13][C3-02][P2][METODO] Tornar a revisão reprodutível e limitar o viés do corpus ISO
% Dar bases/plataformas, strings completas por base, datas, idiomas,
% intervalo temporal, deduplicação, seleção e motivos de exclusão.
% Definir A e B e a regra exata de passagem entre eles. Uma pesquisa que
% exige "ISO 50001/50006" pode omitir trabalhos relevantes de deteção de
% mudança, M&V, modelação energética e fouling que não usam esses termos.
% Justificar o recorte; considerar rastreio de referências e pesquisa
% complementar, declarados como emenda se não estavam no protocolo.
% [/REV-2026-09-13]
\section{Estratégia de pesquisa e seleção}
\label{sec:pesquisa}

%% >> Strings calibradas, screening em tres niveis, extracao assistida por dois modelos independentes sob protocolo auditado (Anexo ref app:validacao).


% [REV-2026-09-13][F04][P2][FIGURA] Fluxo PRISMA com a formação dos dois corpora
% Inserir antes desta secção, no fim da seleção. Mostrar registos por
% fonte, duplicados, triagem, textos completos, exclusões com motivos e
% inclusões; depois a ramificação documentada para Corpus A e B.
% Não preencher contagens estimadas. O fluxo principal deve estar no corpo
% porque define a amostra bibliográfica; o Apêndice B guarda o detalhe.
% Distinguir registos, publicações e estudos quando não coincidam.
% [/REV-2026-09-13]
% [REV-2026-09-13][C3-03][P2][ANALISE] Contagem bibliométrica não demonstra adequação metodológica
% Definir unidade de contagem, denominadores e regras de classificação de
% setor, modelo e validação; admitir múltiplas etiquetas se necessário.
% Um setor pouco representado indica cobertura do corpus, não inexistência
% de métodos nem inadequação dos publicados. Evitar word clouds e redes
% de palavras sem uma pergunta substantiva. Uma tabela por setor/uso/modelo
% e evolução anual com contagens e denominadores costuma ser suficiente.
% [/REV-2026-09-13]
\section{Análise bibliométrica (Corpus A)}
\label{sec:corpus-a}

%% >> Distribuicao da literatura EnPI/EnB por ano, setor e contexto normativo: a evidencia quantitativa do gap.


% [REV-2026-09-13][F05][P2][FIGURA] Mapa de evidência que sustenta a lacuna
% Depois de explicar a extração, usar uma matriz: linhas = estudos ou
% famílias; colunas = contexto industrial, variáveis, horizonte, validação
% temporal, dependência, deriva, cobertura, eventos reais e implantação.
% Codificar "avaliado", "não reportado" e "não aplicável" distintamente,
% com referências nas linhas e texto de síntese por SQ. É a figura/tabela
% mais útil para justificar o que a tua proposta acrescenta. Manter os
% dados de extração completos no Apêndice B; não substituir síntese por cor.
% [/REV-2026-09-13]
\section{Síntese narrativa (Corpus B)}
\label{sec:corpus-b}

%% >> O que os estudos em processo continuo propoem, reconhecem e nao resolvem, por sub-questao.


% [REV-2026-09-13][C3-04][P2][ARGUMENTO] Fecho da revisão tem de ligar evidência e decisões do estudo
% Comparar abordagens existentes, condições em que funcionam, limites de
% validação e o que falta operacionalizar. Terminar com três lacunas
% documentadas e a decisão metodológica que responde a cada uma no Cap. 4.
% Não escrever a conclusão universal prevista no comentário-esqueleto sem
% evidência. Avaliar criticamente pelo menos os estudos mais próximos do
% caso, mesmo que avaliação de risco de viés formal não seja objetivo
% obrigatório de uma scoping review.
% [/REV-2026-09-13]
\section{Síntese: o \emph{gap}}
\label{sec:gap}

%% >> Nao existe metodologia sistematica para avaliar e construir baselines energeticas em processo continuo dentro do enquadramento normativo.


